\section*{Chapitre \ref{ch:g10:drainage-basins} --- Bassins versants, rivières et crues}

\begin{solution}{exo:g10:drainage-basins:1}
Le \omterm{def:g10:drainage-basins:drainage-basin}{bassin versant} d'un point d'une rivière est toute la surface de terres
dont le ruissellement rejoint la rivière en amont de ce point et passe par
lui. La \omterm{def:g10:drainage-basins:drainage-basin}{ligne de partage des eaux} est la ligne, le long des crêtes, qui
sépare deux bassins voisins. Deux gouttes tombées à un mètre l'une de
l'autre, de part et d'autre d'une ligne de partage, descendent des versants
opposés, dans deux bassins différents~: leurs rivières peuvent atteindre
des mers différentes.
\end{solution}

\begin{solution}{exo:g10:drainage-basins:2}
Le \omterm{def:g10:drainage-basins:discharge}{débit} est le volume d'eau qui traverse une section de la rivière chaque
seconde, en \unit{m^3/s}. \qty{1}{m^3/s} représente \qty{1000}{L} chaque
seconde~; en un jour, $1000 \times 86\,400 = 86.4$~millions de litres.
\end{solution}

\begin{solution}{exo:g10:drainage-basins:3}
Avant la crue, la rivière transportait environ \qty{330}{m^3/s} (entre
environ \num{240} et \qty{390}{m^3/s}). Le pic a été de \qty{9231}{m^3/s}
le 21~janvier. La montée a commencé le 18~janvier (\qty{362}{m^3/s})~: elle
a duré trois jours.
\end{solution}

\begin{solution}{exo:g10:drainage-basins:4}
$1/50 = 0.02$, soit \qty{2}{\percent}. Non~: les années sont indépendantes~;
la probabilité est de nouveau de \qty{2}{\percent} cette année, quoi qu'il
se soit passé l'an dernier.
\end{solution}

\begin{solution}{exo:g10:drainage-basins:5}
Un pied cube vaut $0.3048^3 \approx \qty{0.02832}{m^3}$, donc
$593\,000 \times 0.02832 \approx \qty{16800}{m^3/s}$. L'Amazone transporte
$206\,000/16\,800 \approx 12$ fois plus.
\end{solution}

\begin{solution}{exo:g10:drainage-basins:6}
Aires des bandes~: $1.5 \times 0.3 = 0.45$, $0.90$, $1.05$ et \qty{0.60}{m^2}~:
$A = \qty{3.0}{m^2}$. \omterm{def:g10:drainage-basins:discharge}{Débits}~: $0.2 \times 0.45 = 0.09$, $0.5 \times 0.90
= 0.45$, $0.6 \times 1.05 = 0.63$, $0.3 \times 0.60 = 0.18$~: $Q =
\qty{1.35}{m^3/s}$. Vitesse moyenne~: $1.35/3.0 = \qty{0.45}{m/s}$.
\end{solution}

\begin{solution}{exo:g10:drainage-basins:7}
$66 \times 10^6 \times 0.020 = \qty{1.32e6}{m^3}$. Deux jours font
$2 \times 86\,400 = \qty{172800}{s}$~: $1.32 \times 10^6/172\,800 \approx
\qty{7.6}{m^3/s}$ en moyenne.
\end{solution}

\begin{solution}{exo:g10:drainage-basins:8}
Classés~: 120, 88, 65, 59, 51, 47, 42, 38, 35. Avec $n = 9$, $T = 10/m$~: le
plus grand (120) a $T = 10$~ans, le deuxième (88) 5~ans, le troisième (65)
environ 3.3~ans. La médiane, 51, a le rang 5 et $T = 2$~ans~: elle est
égalée ou dépassée une année sur deux.
\end{solution}

\begin{solution}{exo:g10:drainage-basins:9}
L'\omterm{def:g10:drainage-basins:hydrograph}{hydrogramme} de la ville monte plus tôt et bien plus haut, et redescend
vite~: sur des pentes fortes, des toits et des routes, peu d'eau s'infiltre
et le ruissellement atteint le ruisseau vite et d'un seul coup~; son \omterm{def:g10:drainage-basins:hydrograph}{temps
de réponse} est court. L'\omterm{def:g10:drainage-basins:hydrograph}{hydrogramme} du bassin boisé monte plus tard, vers un
pic plus bas et plus large~: le sol et la forêt absorbent une grande partie
de la pluie et la relâchent lentement, et l'eau s'écoule lentement sur des
pentes douces. Son \omterm{def:g10:drainage-basins:hydrograph}{débit de base} reste plus élevé après l'orage.
\end{solution}

\begin{solution}{exo:g10:drainage-basins:10}
$9231 \times 86\,400 \approx \qty{7.98e8}{m^3}$, répartis sur
$29\,900 \times 10^6 = \qty{2.99e10}{m^2}$~: environ \qty{0.027}{m}, soit
\qty{27}{mm}. Le \omterm{def:g10:drainage-basins:hydrograph}{débit de base} donne $330 \times 86\,400/2.99 \times 10^{10}
\approx \qty{0.001}{m}$, environ \qty{1}{mm} par jour~: le jour du pic a
transporté environ 28 fois plus.
\end{solution}

\begin{solution}{exo:g10:drainage-basins:11}
$P = 1 - 0.9^{10} \approx 1 - 0.349 = 0.651$~: environ \qty{65}{\percent}.
Elle ne survient pas du tout avec une probabilité $0.9^{10} \approx 0.35$,
environ \qty{35}{\percent}.
\end{solution}

\begin{solution}{exo:g10:drainage-basins:12}
L'\omterm{def:g10:drainage-basins:hydrograph}{hydrogramme} est un triangle de trois jours de long et de
\qty{1200}{m^3/s} de haut~: son aire vaut $\tfrac12 \times 1200 \times 3 \times 86\,400 \approx
\qty{1.56e8}{m^3}$, soit \num{156}~millions de mètres cubes. La partie
au-dessus de \qty{800}{m^3/s} est un triangle plus petit de même forme,
trois fois plus petit en hauteur ($400$ au lieu de $1200$) et en durée (un
jour au lieu de trois)~: $\tfrac12 \times 400 \times 86\,400 \approx \qty{1.73e7}{m^3}$,
soit \num{17}~millions de mètres cubes. C'est moins que les
\num{60}~millions que peut stocker la retenue~: oui.
\end{solution}

\begin{solution}{exo:g10:drainage-basins:13}
Pleine jusqu'à ses berges, la rivière est profonde de $800/(100 \times 2) = \qty{4}{m}$.
(1) La plaine d'inondation fait passer \qty{400}{m^3/s} sur \qty{900}{m} à
\qty{0.5}{m/s}~: sa profondeur vaut $400/(900 \times 0.5) \approx \qty{0.9}{m}$.
(2) Entre les digues~: $1200/(100 \times 2) = \qty{6}{m}$ de profondeur,
\qty{2}{m} au-dessus de ses berges, plus de deux fois la profondeur sur la
plaine sans digues.
\end{solution}

\begin{solution}{exo:g10:drainage-basins:14}
On cherche $(1 - 1/T)^{30} = 0.5$, donc $1 - 1/T = 0.5^{1/30} \approx 0.97716$,
$1/T \approx 0.02284$ et $T \approx 44$~ans.
\end{solution}

\begin{solution}{exo:g10:drainage-basins:15}
Chaque carré envoie son eau vers son voisin le plus bas (lignes du nord au
sud, colonnes de l'ouest à l'est)~:
\begin{center}
\begin{tabular}{|c|c|c|c|}
\hline
exutoire & ouest & ouest & ouest \\ \hline
nord & ouest & ouest & sud \\ \hline
nord & nord & sud & sud \\ \hline
nord & est & est & exutoire \\ \hline
\end{tabular}
\end{center}
En suivant les flèches, les carrés 47, 39, 44, 37, 46 et 30 s'écoulent vers
l'exutoire sud-est~: 6~carrés, \qty{0.06}{km^2}. Les 10 autres carrés
s'écoulent vers l'exutoire nord-ouest~: \qty{0.10}{km^2}. La ligne de
partage part du bord est entre 50 et 47, puis passe entre 48 et 47, entre
48 et 44, entre 49 et 44, entre 49 et 46, et entre 51 et 46 jusqu'au bord
sud.
\end{solution}

\begin{solution}{pb:g10:drainage-basins:1}
\textbf{1.} $11\,560 \times 1.609^2 \approx \qty{29900}{km^2}$.

\textbf{2.} La \omterm{def:g10:drainage-basins:drainage-basin}{ligne de partage des eaux} du bassin situé en amont de la
station.

\textbf{3.} $484\,000 \times 0.3048^3 \approx \qty{13700}{m^3/s}$.

\textbf{4.} $13\,705/3341 \approx 4.1$ fois.

\textbf{5.} $A = 400 \times 8 = \qty{3200}{m^2}$ et $v = Q/A = 13\,705/3200
\approx \qty{4.3}{m/s}$.

\textbf{6.} Environ \qty{330}{m^3/s}.

\textbf{7.} \qty{9231}{m^3/s}, le 21~janvier.

\textbf{8.} Trois jours.

\textbf{9.} $9231 \times 86\,400/(29\,900 \times 10^6) \approx
\qty{0.027}{m}$~: environ \qty{27}{mm}.

\textbf{10.} De l'eau nouvelle a atteint la rivière~: une autre chute de
pluie sur le bassin (ou une fonte de neige) a envoyé une seconde impulsion
de ruissellement, plus faible.

\textbf{11.} Dix-sept jours après le pic. La montée vient du ruissellement
des pentes proches des ruisseaux, qui arrive vite et d'un seul coup~; sur la
courbe de décrue, la rivière reçoit encore l'eau des parties les plus
éloignées du bassin et l'eau qui s'écoule lentement hors du sol.

\textbf{12.} $T = n + 1 = 96$~ans.

\textbf{13.} $96/2 = 48$~ans.

\textbf{14.} $96/3 = 32$~ans.

\textbf{15.} $96/4 = 24$~ans et $96/5 = 19.2$~ans. Peu~: dans les deux cas,
ce \omterm{def:g10:drainage-basins:discharge}{débit} revient environ une fois en vingt ans.

\textbf{16.} Cinq années (1936, 1937, 1943, 1972 et 1996)~: $5/95 \approx
0.053$, environ \qty{5}{\percent} par an, soit une \omterm{def:g10:drainage-basins:return-period}{période de retour}
d'environ $96/5 \approx 19$~ans.

\textbf{17.} $1/100 = 0.01$, soit \qty{1}{\percent}.

\textbf{18.} $1 - 0.01 = 0.99$.

\textbf{19.} $0.99^{30} \approx 0.740$.

\textbf{20.} $1 - 0.740 = 0.26$~: \textbf{la probabilité d'au moins une crue
centennale en 30~ans est d'environ \qty{26}{\percent}}, une sur quatre.
\end{solution}
